\documentclass[11pt]{article}

% Page layout and typography
\usepackage[margin=1in]{geometry}
\usepackage{microtype}
\usepackage{parskip}

% Mathematics
\usepackage{amsmath,amssymb,amsthm,mathtools}
\usepackage{bm}

% Figures and links
\usepackage{graphicx}
\graphicspath{{figures/}}
\usepackage{xcolor}
\usepackage[hidelinks]{hyperref}

% Lists
\usepackage{enumitem}
\setlist{nosep}

% Theorem environments
\theoremstyle{definition}
\newtheorem{definition}{Definition}[section]
\newtheorem{example}[definition]{Example}
\newtheorem{remark}[definition]{Remark}

\theoremstyle{plain}
\newtheorem{theorem}[definition]{Theorem}
\newtheorem{lemma}[definition]{Lemma}
\newtheorem{proposition}[definition]{Proposition}
\newtheorem{corollary}[definition]{Corollary}

% Common number systems
\newcommand{\R}{\mathbb{R}}
\newcommand{\C}{\mathbb{C}}
\newcommand{\Q}{\mathbb{Q}}
\newcommand{\Z}{\mathbb{Z}}
\newcommand{\N}{\mathbb{N}}
\newcommand{\D}{\mathbb{D}}
\newcommand{\T}{\mathbb{T}}
\newcommand{\RS}{\mathbb{S}^2}
\newcommand{\per}{\text{per}}

% Common operators
\DeclareMathOperator{\im}{im}
\DeclareMathOperator{\rank}{rank}
\DeclareMathOperator{\Span}{span}
\DeclareMathOperator{\id}{id}


\title{MATH4341 Lecture Notes}
\author{Caelen Gelineau}
\date{\today}

\begin{document}

\maketitle
\tableofcontents

\newpage
\section{Lecture 9 - 23/09/26}
\begin{definition}
    A \textit{subbasis} $\mathcal{C}$ for a topology on $X$ is a collection of subsets that cover $X$, meaning that their union is all of $X$.
\end{definition}
\begin{example}
    Given $f: \R \rightarrow \R^\N = \R^\omega$, $f(t) = (t,t,\dots, t, \dots)$ and each $f_i: \R \rightarrow \R$, we want 
    \[
        f \text{ continuous} \iff f_i \text{ continuous}
    \]
   in other words,
   \[
    f: Y \rightarrow \prod_{i\in I}^{\infty} X_i \text{ continuous} \iff \text{each } f_i: Y \rightarrow X_i \text{ continous} 
   \]
    Claim: $f: \R \rightarrow (\R^\omega, \text{box topology})$ is \textbf{NOT} continuous.
    Take basis element $B = \prod_{n = 1}^{\infty} (-\frac{1}{n}, \frac{1}{n})$ and we claim $f^{-1}(B)$ is not open in $\R$.
    Note that $x = 0 \in f^{-1}(B)$ since $f(0) = \overrightarrow{0} \in B$. Assume it is open. Then we have an open interval $\prod_{n=1}^{\infty}(-\frac{1}{\epsilon}, \frac{1}{\epsilon}) \in B $ around $x = 0$.
    But this implies for the box topology $(-\frac{1}{\epsilon}, \frac{1}{\epsilon}) \in (-\frac{1}{n}, \frac{1}{n}) \forall n$ which is not true, so we need a new topology.\\
    Since $f_i = \pi_i \circ f$, we need $\pi_i$ to be continuous to satisfy the desired condition. So find a topology for which this is true. We want
    \[
        \pi_j: \prod_{i \in I}^{\infty} X_i \rightarrow X_j 
    \]
    to be continuous, so find a topology to put on the product space.
    \[
        \pi_j^{-1} \text{ is open in } \prod_{i \in I}^{\infty} X_i  
    \]
    notice that when projecting onto an open set along an axis for $\R^2$, we have
    \[
        \pi_j^{-1} = U_1 \times X_2
    \]
    for an open set $U_1$ in $X_j$. Thus the preimage is a strip where one dimension is unrestricted. The same applies for every other dimension. 
    However, these strips are not a basis because the infinite vertical strips do not satisfy (B2) because the intersection of two strips cannot contain an infinite strip.
    Thus we create a subbasis so that our finite rectangles are subbasis elements.
    \[
        \pi_j^{-1}(U_j) = \prod_{i \in I}^{\infty} V_i,  
    \]
    where
    \[
        V_i =
        \begin{cases}
            X_i & \text{if } i \ne j \\
            U_j & \text{if } i = j
        \end{cases}
    \]
    Also note that a basis $=$ finite intersection of a subbasis $= \prod_{i \in I}^{\infty} V_i$ where $V_i = X_i$ for all but finitely many $i$'s.
    This is the product topology.
    \[
       subbasis \rightarrow_{\text{finite intersection}}basis \rightarrow_{\text{union}} topology
    \]
    The subbasis allows us to include the missing scalar multiplication operation for vector spaces that we are missing when we begin from just a basis.
\end{example}
\begin{theorem}
   \[
        f: Y \rightarrow \prod_{i\in I}^{\infty} X_i \text{ continuous} \iff \text{each } f_i: Y \rightarrow X_i \text{ continous} 
   \]
    
\end{theorem}
\begin{proof}
    ($\Rightarrow$) Assume $f$ is continuous. We want to show $f^{-1}(U)$ is open in $Y$. It is enough to consider $U$ as a subbasis because subbasis implies basis holds and topology holds.
    The preimage of the projection is of the form
    \[
        \pi_j^{-1}(U_j) = \prod_{i \in I}^{\infty} U_i 
    \]
    where $U_i = X_i\text{ if } i \ne j$.
so then a subbasis element is
\[
    f^{-1}(\pi_j^{-1}(U_j)) = (\pi_j \circ f)^{-1}(U_j) = f_j^{-1}(U_j) = f_j
\]
thus since the subbasis elements are equal to $f_j$ the preimage of each are open in $Y$ thus $f_j$ is continuous.

$f^{-1}(U)$ open in $Y$ for subbasis $U$ implies the basis $B = U_1 \cap \dots \cap U_n$ which implies $f^{-1}(B) = f^{-1}(U_1)\cap \dots$ implies open which implies holds true for topology.\footnote{Show this at home. Also show that this proof fails for box topology and why we need product topology (to show $\pi_i$ continuous for forward direction)}\\
($\Leftarrow$) Trivial, in slides.

\end{proof}
\end{document}